\paragraph{What runtimes should do}
The rules of Section~\ref{sec:theory} cost nothing to apply. A runtime that streams layers with prefetching should spread its resident layers evenly instead of filling GPU memory in model order, and should size its ring by the ratio of copy to compute time per layer: two slots when a copy takes at least twice as long as a layer's computation, three otherwise, or the faster of the two according to the evaluator, which needs only per-layer compute and copy times that a runtime can measure in one iteration. When a layer computes longer than it takes to copy, resident layers do not shorten the iteration at all, and the memory is better spent on a larger batch or key--value cache. When copies dominate even with every free byte resident, placement does not matter either, and only fewer bytes (quantisation) or a faster path help.

\paragraph{Faster paths and CXL memory}
The analysis depends on the path only through the copy time per layer $f$. A CXL Type-3 expander reached through a page-locked host mapping is such a path \cite{yoon2025tract,jang2026itme}, and reported bandwidths for current devices are of the order we measured on PCIe Gen3. If copies from such a device take a constant time per layer, the theorems apply to it as stated. Whether they do is an open question that our measurements cannot answer: copies from an unpinned CXL region go through a bounce buffer \cite{yoon2025tract}, and contention on a shared device or with host memory traffic would make $f$ vary from copy to copy. A faster path, such as several expanders interleaved or PCIe Gen5, lowers $f$ and moves a given workload towards the regime $f\le c$ in which placement no longer matters; a GPU that computes faster moves it the other way.

\paragraph{Power-limited GPUs}
On the power-limited T4 a concurrent copy and long runs lowered the SM clock, which changes $c$ during an iteration and breaks assumption A4. The effect was large enough to reverse the order of two placements in one configuration of our first runs (Section~\ref{sec:res-power}). The contention variant of the model captures the part caused by concurrent copies but not the slower drift of the clock with temperature and power. Data-centre GPUs with higher power limits are less affected, but any deployment that runs close to its power limit should measure $c$ under the same conditions as the iteration it wants to predict.

\paragraph{Limitations}
(i) All measurements are on one GPU type, a power-limited T4 on Colab virtual machines, with two model shapes. The clock lock that makes the measurements reproducible runs the GPU at about a third of its maximum clock, which shifts every workload towards larger $f/c$; on a faster GPU or at default clocks the same workload has a smaller $c$ and the regime boundaries move. (ii) The layer kernels omit rotary embeddings, expand grouped key--value heads explicitly, and share one read-only key--value cache across layers; each configuration is measured once, with one random placement per cell. (iii) The runtime is written in Python without CUDA graphs, and in the shortest decode iterations host launch overhead, which the model ignores, may be on the critical path. (iv) We stream whole layers; a runtime that splits layers into tensors has more, smaller units, to which Lemma~\ref{lem:cycle} applies unchanged but for which we did not measure. (v) The key--value cache stays in GPU memory; streaming it as well adds per-layer copy time that grows with the context, which fits the model as a larger $f_i$ but was not measured. (vi) We assume one copy engine and in-order execution; runtimes that reorder layers or split copies across several engines need a different event graph. (vii) The theorems cover layers of equal cost and require $k\le m$; for non-uniform units we offer an exact evaluator and a search, not a closed form.
